\section{Décima iteración de la Fase 3: Unir segmentos y dígitos en pantalla}
\label{sec:UnirSegmentos}

Esta iteración rehace el mecanismo de unión de segmentos de arriba abajo y
convierte su criterio en la regla general de todos los agrupamientos. Con esto resuelto, se puede finalziar el desarrollo de incluir contadores en la pantalla.

\subsection{Nuevo sistema de unión}

La unión \emph{fusionaba} los segmentos marcados en uno de los existentes: los
absorbidos desaparecían y el superviviente heredaba sus píxeles, sus enlaces
re-apuntados y un número de mapa heredado también. El sistema cambia para crear un nuevo segmento para la unión, y a partir de qauí ya es sencillos de ajustar para adaptar sus propiedades.

\begin{itemize}
	\item \textbf{Unir crea una pieza nueva}. Al unir, por ejemplo, los segmentos
	\texttt{seg\_020} y \texttt{seg\_030}, no se fusionan en uno de los dos: se crea un
	\emph{segmento} nuevo que se llama (\texttt{seg\_F01}), y una \emph{imagen} nueva (la fusión de ambas), que llamaremos \texttt{img\_F01}. Los
	segmentos e imágenes originales \textbf{siguen existiendo}, marcados como
	\emph{fusionados con F01}: no se presentan en pantalla ni se puden manejar, pero sus datos 
	siguen en el cartucho y la autoría los conserva. 
	
	\item \textbf{Un único sistema de absorción}. «Unir», «Hacer dígito» y
	«Hacer contador» son el mismo mecanismo: crear una pieza nueva (con
	un número Fnn correlativo) que absorbe a los originales, hereda sus
	propiedades (el primer valor no nulo de cada campo) y se le re-apuntan
	todos los enlaces externos que les referenciaban. «Separar» es siempre la
	operación inversa, también deshacible.
	
	\item \textbf{Los fusionados no se presentan}. Un segmento fusionado lleva
	la marca \texttt{visible:false} en el cartucho: nadie debe encenderlo ni
	pintarlo (su imagen está dentro de la fusión de su pieza). Al generar el
	cartucho para el ESP32 no pasarán a este los segmentos invisibles ni sus imágenes.
\end{itemize}

Sobre esta base, los contadores que no fue posible agrupar correctamente en la iteración anterior, se pueden ya montar y presentar sin problemas en la pantalla.

\subsection{El dígito de siete segmentos}
\label{subsec:SieteSegmentosF3I10}

El autor del cartucho dibuja un «8» en cada posición del número y el sistema sabe
qué barras enciende cada dígito (un 1 es \texttt{b+c}, un 2 es
\texttt{a,b,g,e,d}, etcétera): la tabla es fija y la trae el sistema. Con esto
el autor no dibuja los diez dígitos posibles ni declara qué número representa
cada uno. El 8 de base se puede dibujar tanto normal como en cursiva, el único requisito
es que se puedan extraer correctamente los 7 segmentos del gráfico unívocamente.

Los elementos que se definen son:
\begin{itemize}
	\item \textbf{Dígito}: el autor selecciona recuadrando con el ratón las siete barras de un
	«8» y pulsa \emph{«Hacer dígito»} (o simplemente \emph{«Unir»}: con siete
	barras de un «8» reconocible, la unión \emph{es} hacer el dígito): el
	sistema clasifica cada barra por su \textbf{geometría} dentro del 8 (las
	letras \texttt{a}--\texttt{g} se asignan solas: las tres horizontales por
	altura son \texttt{a}, \texttt{g}, \texttt{d}; las verticales de la
	derecha \texttt{b} y \texttt{c}, las de la izquierda \texttt{f} y
	\texttt{e}) y crea un segmento \textbf{NUEVO} de tipo dígito, llamado
	como su imagen fusionada (\texttt{F01}, \texttt{F02}, \dots), que
	\textbf{agrupa las barras} mediante los enlaces \texttt{barra\_a} a
	\texttt{barra\_g} (siete enlaces; el dígito no \emph{es} ninguna barra, las
	\emph{agrupa}). El dígito nuevo lleva una \textbf{nueva imagen que fusiona
		las siete barras} (mapa de bits OR de sus píxeles, el mismo criterio que
	\emph{Unir}): a partir de ese momento el editor maneja el dígito como
	una sola pieza con esa imagen, y \textbf{las siete barras individuales dejan de verse}
	en el mapa y en la pantalla de autoria (clicarlas selecciona el dígito),
	pero \emph{siguen viajando en el cartucho} con sus sprites: el motor
	necesita encenderlas una a una.
	Como en \emph{Unir}, el grupo es una pieza nueva y \emph{Separar} la
	elimina devolviendo las barras al editor; todo en una acción atómica del
	historial, deshacible de un gesto. Un dígito \textbf{no se une con nada}
	ni puede cambiarse de tipo, ya que la pieza del «8» es indivisible, y si se cambia el tipo esta información (las siete barras que lo componen) se perdería.
	
	\item \textbf{Contador}: el autor marca los $n$ dígitos que forman el número (\textbf{limitado a 4 dígitos como máximo}) y
	pulsa \emph{«Hacer contador»} (con \emph{un solo} dígito seleccionado, su
	propio panel también lo ofrece: un número de una sola posición es un
	contador legítimo): se crea otro segmento \textbf{nuevo} de tipo
	contador, llamado como su imagen (\texttt{Fxx}), que
	\textbf{agrupa los dígitos} mediante los enlaces \texttt{digito\_1} a
	\texttt{digito\_4}: \texttt{digito\_1} es el dígito de las
	\textbf{unidades} (el más a la derecha; es \emph{obligatorio}) y los
	demás crecen hacia la izquierda. A partir de entonces solo se ve y se
	maneja el contador completo con todos sus dígitos: \textbf{los dígitos individuales dejan de ser accesibles} en el editor
	(clicar en cualquier número selecciona el contador completo), aunque siguen viajando individualmente en el
	cartucho: el motor enciende las barras reales de cada posición. Como el
	dígito, un contador \textbf{no se puede unir con nada} ni cambia su tipo. Sus
	propiedades son:
	\begin{itemize}
		\item \textbf{Contador}: Indica cual de los 4 contadores se asocia al número.
		Se \textbf{autoasigna} al crear el contador al primero que no use
		otro contador; el panel permite cambiarlo (no se aceptan contadores sin
		número: la verificación exige 1--4).
		\item \textbf{Valor inicial}: Por defecto sera 0, pero en caso de los contadores de vidas, se debe establecer el valor inicial adecuado.
		\item \textbf{Tipo de cuenta}: si es de cuenta ascendente o descendente. La ascendente es idenfinida, la descendente hasta cero.
		\item \textbf{Presentar el cero}: Indica si el valor \textbf{cero} se debe presentar o dejar la zona completamente en blanco.
		\item \textbf{Cuántos dígitos presenta}: se cuentan automáticamente los
		dígitos agrupados (\texttt{digito\_1} a \texttt{digito\_4}) y se muestra
		como \emph{solo lectura}. Sólo sirve como verificación de que se agrupó
		lo que el autor quería.
	\end{itemize}
	
	\item \textbf{Separar}: el panel del dígito y del contador ofrecen
	\emph{«Separar»}, que disuelve el grupo y devuelve sus piezas al editor
	(acción inversa, deshacible con todo lo demás).
	
	\item \textbf{El encendido}: en cada refresco el contador toma el valor del
	contador interno asociado, lo parte en dígitos y enciende en cada posición
	(el dígito real que el autor dibujó, con sus barras) las que la tabla marca
	para cada dígito; el resto apagadas. Una posición en blanco es simplemente
	las siete apagadas. Un contador descendente no baja de cero: al llegar a 0
	se queda presentándolo.
\end{itemize}

\subsection{Decisiones menores}
\label{subsec:DecisionesF3I10}

\begin{itemize}
	\item \textbf{Segmentos ficticios}: dígito y contador son piezas
	\emph{ficticias} creadas al agrupar, igual que la pieza de una unión.
	Se llaman como su imagen (\texttt{F01}, \texttt{F02}, \dots: la pieza
	lleva el nombre de su sprite, no un correlativo \texttt{seg\_x}), no
	desplazan los índices de los segmentos reales y desaparecen al
	\emph{Separar}. En el GWI el ficticio viaja con su imagen fusionada como
	sprite, y además exporta sus enlaces: el dígito las siete barras
	(\texttt{barra\_a} a \texttt{barra\_g}) y el contador sus dígitos
	(\texttt{digito\_1} a \texttt{digito\_4}) más
	\texttt{contadorInterno}, \texttt{valorInicial},
	\texttt{direccion} y \texttt{presentarCero} (solo si es falso). Los
	dígitos no se comparten entre contadores y las barras no se comparten
	entre dígitos (verificación). Separar un dígito que está dentro de un
	contador exige separar primero el contador.
	
	\item \textbf{Visibilidad declarada}: todo segmento lleva la propiedad
	\texttt{visible}. La calcula el diseñador al reconstruir (un segmento
	absorbido o agrupado está \emph{dentro} de la imagen fusionada de su
	pieza: \texttt{visible:false}) y viaja en el GWI, de modo que ni el
	motor ni la previsualización recalculan nada: quien lleva
	\texttt{visible:false} no se maneja, enciende ni se pinta nunca. Es la
	preparación del generador del cartucho para el ESP32: los fusionados
	viajan hoy marcados (sus imágenes siguen en el paquete) y allí se
	eliminarán directamente.
	
	\item \textbf{Los contadores internos son genéricos}: Hasta cuatro por cartucho que llamaremos \textbf{contador 1, contador 2, contador 3 y contador 4}. Son
	independientes de los segmentos de tipo vida/muerte, que son 
	iconos que los representan (p. ej. en el parachute los 3 tiburones para las muertes).
	Ambos sistemas conviven, y el manejo de los segmentos
	vida/muerte se definirán en otra iteración, junto al evento que determine
	cuándo termina la partida.
	
	\item El contador \textbf{interno} sigue creciendo independientemente de que
	la presentación se quede sin dígitos, y por debajo de cero no baja: el fin
	de partida necesitará el valor real.
	
	\item Relleno de la izquierda: los valores menores que
	$10^n$ van rellenos con \emph{blancos} a la izquierda si lo necesitan
	($\sqcup$0, $\sqcup$1, \dots, 10, 11, \dots), los valores mayores o iguales a $10^n$ se
	presentan rellenos con ceros a la izquierda (\dots, 98, 99, 00, 01, \dots).
	
	\item Las propiedades \texttt{ColisionPositiva}/\texttt{ColisionNegativa} de los
	desenlaces eligen cuál de los cuatro contadores aumenta (o ninguno); la
	decisión sigue tocando el contador elegido una sola vez.
\end{itemize}

\subsection{«Posicionar juntos» retirado}

La reposición en pantalla de los segmentos definida en \ref{sec:Posicionar} (presentar un segmento en las mismas
coordenadas de otro) nació para presentar el contador numéricos de las vidas, definiendo cada dígito como una imagen que se presentaría en el mismo lugar de pantalla. Esto era menos visual ya que obligaba a tener todos los dígitos separados, aunque luego se presentasen juntos, ocupaban una zona en la pantalla de diseño, que luego desaparecía.

Con los contadores reales ya no tiene sentido y se retira esta funcionalidad
del editor; la maquinaria interna desaparece pues no se necesita.

\begin{figure}[H]
	\centering
	\includegraphics[width=0.7\linewidth]{imagenes/propias/c05-I013.png}
	\caption{Numeros posicionables}
	\label{fig:c05-I013_B.png}
\end{figure}

\clearpage